\section{The Four Happinesses}

Taoism holds to the ideal of the four happinesses. The first two concern the extension of the individuality and are included in the restoration of the primordial state.

\begin{itemize}

\item \textbf{Longevity}: The perpetuity of individual existence.

\item \textbf{Posterity}: The indefinite prolongations of the individual through all his modalities.
\end{itemize}


The remaining two refer to the higher and extra-individual states of the being.

\begin{itemize}
\item \textbf{Great Wisdom:} Learned master.

\item \textbf{Perfect Solitude:} Peace of mind, inner quietude
\end{itemize}

The fifth happiness contains them all principally and unites them synthetically, i.e., the \textbf{Supreme Identity}.

From this, we see that longevity for its own sake is of no value. In popular culture, we admire elders who manage to accomplish physical tasks more appropriate to younger men, while precious few are interested in the state of Perfect Solitude.

Longevity, extending beyond one's reproductive capacity, makes no biological sense. Hence, the value of longevity is to provide more time to develop one's Real Self and to attain merits.


\flrightit{Posted on 2019-10-05 by Losang Shenphen}

\begin{center}* * *\end{center}

\begin{footnotesize}
\begin{sffamily}
\texttt{AghorNath on 2019-10-05 at 11:24 said:}

God bless you brother! Praying for your healing.

\hfill

\texttt{Sviatoslav on 2019-10-06 at 07:56 said:}

His work is so unappreciated by today's world. Only a few know about Calogero's heroic mission. God bless him.


\hfill

\end{sffamily}
\end{footnotesize}

